\chapter{Mecanismos de reacción: flechas curvas}\label{ch:g12:curly-arrows}

Una \omterm{def:g8:balanced-equations:equation}{ecuación ajustada} es como la primera y la última fotografía de una
película: muestra los \omterm{def:g7:chemical-reactions:reactant}{reactivos} antes y los productos después, y nada en
medio. Sin embargo, las \omterm{def:g7:atoms-and-molecules:molecule}{moléculas} no se limitan a intercambiar \omterm{def:g7:atoms-and-molecules:atom}{átomos}.
Los enlaces se rompen y se forman de uno en uno, a medida que pares de
\omterm{def:g9:inside-the-atom:nucleus}{electrones} pasan de un lugar a otro, a través de especies de vida corta
que nunca aparecen en la ecuación. Los químicos dibujan estos movimientos
con \omterm{def:g12:curly-arrows:curly-arrow}{flechas curvas}. Con unas pocas reglas, estas flechas explican por qué
una reacción ocurre donde ocurre, y predicen reacciones nunca vistas.

\begin{recall}
La \omterm{def:g11:polarity-and-cohesion:electronegativity}{electronegatividad}, los \omterm{def:g11:polarity-and-cohesion:polar-bond}{enlaces polares} y las \omterm{def:g11:polarity-and-cohesion:polar-bond}{cargas parciales}
(\cref{ch:g11:polarity-and-cohesion}). Las \omterm{def:g10:lewis-and-shape:lewis-structure}{estructuras de Lewis}: \omterm{def:g10:lewis-and-shape:lone-pair}{pares
enlazantes} y \omterm{def:g10:lewis-and-shape:lone-pair}{pares libres} (\cref{ch:g10:lewis-and-shape}). Los \omterm{def:g11:functional-groups:functional-group}{grupos
funcionales} y las familias (\cref{ch:g11:functional-groups}).
\end{recall}

\section{Sitios ricos y pobres en electrones}

\begin{definition}[Sitios dadores y aceptores de electrones]\label{def:g12:curly-arrows:site}
Un \emph{sitio dador de electrones}\index{sitio dador de electrones} de una
\omterm{def:g7:atoms-and-molecules:molecule}{molécula} o de un \omterm{def:g9:ions:ion}{ion} es un lugar rico en \omterm{def:g9:inside-the-atom:nucleus}{electrones}, capaz de ceder un
par para formar un nuevo enlace: un \omterm{def:g10:lewis-and-shape:lone-pair}{par libre}, un \omterm{def:g10:lewis-and-shape:multiple-bond}{doble enlace}, un \omterm{def:g7:atoms-and-molecules:atom}{átomo}
con carga negativa o con \omterm{def:g11:polarity-and-cohesion:polar-bond}{carga parcial} $\delta^-$. Un \emph{sitio aceptor
de electrones}\index{sitio aceptor de electrones} es un lugar pobre en
\omterm{def:g9:inside-the-atom:nucleus}{electrones}, capaz de recibir un par: un \omterm{def:g7:atoms-and-molecules:atom}{átomo} con carga positiva o con
\omterm{def:g11:polarity-and-cohesion:polar-bond}{carga parcial} $\delta^+$.
\end{definition}

\begin{example}[Los sitios de dos moléculas]\label{ex:g12:curly-arrows:sites}
En el bromoetano, \ce{CH3-CH2-Br}, el bromo es más electronegativo que el
carbono: el carbono unido a él lleva $\delta^+$ (sitio aceptor), y el
bromo, $\delta^-$ y tres \omterm{def:g10:lewis-and-shape:lone-pair}{pares libres}. En la propanona, el oxígeno del
grupo \ce{C=O} lleva $\delta^-$ y dos \omterm{def:g10:lewis-and-shape:lone-pair}{pares libres} (sitios dadores), y su
carbono, $\delta^+$ (sitio aceptor). El \omterm{def:g9:ions:ion}{ion} hidróxido \ce{OH-}, con su
carga negativa y sus tres \omterm{def:g10:lewis-and-shape:lone-pair}{pares libres}, es un dador fuerte.
\end{example}

\begin{omfigure}
\setchemfig{atom sep=2.2em}
\begin{tabular}{cc}
\chemfig{H_3C-C(-[2]H)(-[6]H)-\charge{90:2pt=\:,0=\:,270:2pt=\:}{Br}} &
\chemfig{H_3C-C(-[6]CH_3)=[2]\charge{0=\:,180=\:}{O}} \\[2pt]
\small $\delta^+$ en el carbono unido al \ce{Br}, $\delta^-$ en el \ce{Br} &
\small $\delta^+$ en el carbono del \ce{C=O}, $\delta^-$ en el \ce{O} \\
\small bromoetano & \small propanona
\end{tabular}

{\small \omterm{def:g7:atoms-and-molecules:atom}{Átomos} de carbono pobres en \omterm{def:g9:inside-the-atom:nucleus}{electrones} ($\delta^+$, sitios
aceptores) junto a \omterm{def:g7:atoms-and-molecules:atom}{átomos} electronegativos ricos en \omterm{def:g9:inside-the-atom:nucleus}{electrones}
($\delta^-$, con \omterm{def:g10:lewis-and-shape:lone-pair}{pares libres}: sitios dadores).}
\end{omfigure}

\section{Flechas curvas}

\begin{definition}[Flecha curva]\label{def:g12:curly-arrows:curly-arrow}
Una \emph{flecha curva}\index{flecha curva} representa el movimiento de un
par de \omterm{def:g9:inside-the-atom:nucleus}{electrones} durante una etapa de una reacción. Sale del par que se
desplaza, un \omterm{def:g10:lewis-and-shape:lone-pair}{par libre} o un enlace, y termina donde va el par: en un
\omterm{def:g7:atoms-and-molecules:atom}{átomo}, donde forma un nuevo enlace, o en un \omterm{def:g7:atoms-and-molecules:atom}{átomo} que se queda con el par
de un enlace que se rompe.
\end{definition}

\begin{method}[Dibujar flechas curvas]\label{met:g12:curly-arrows:drawing}
\begin{enumerate}
  \item Localiza el sitio dador y el sitio aceptor.
  \item Dibuja una flecha desde el par dador (un \omterm{def:g10:lewis-and-shape:lone-pair}{par libre} o un enlace)
        hasta el \omterm{def:g7:atoms-and-molecules:atom}{átomo} aceptor: allí se forma un nuevo enlace.
  \item Si el \omterm{def:g7:atoms-and-molecules:atom}{átomo} aceptor quedara entonces con demasiados \omterm{def:g9:inside-the-atom:nucleus}{electrones}
        (más de un octeto, o más de dos en el caso del hidrógeno), dibuja
        una segunda flecha que rompa uno de sus enlaces, con el par
        yendo hacia el \omterm{def:g7:atoms-and-molecules:atom}{átomo} más electronegativo.
  \item Comprueba las cargas tras la etapa: el \omterm{def:g7:atoms-and-molecules:atom}{átomo} que cedió un \omterm{def:g10:lewis-and-shape:lone-pair}{par
        libre} gana una unidad de carga positiva, y el que recibió un par
        como \omterm{def:g10:lewis-and-shape:lone-pair}{par libre}, una de carga negativa; la carga total no cambia.
\end{enumerate}
\end{method}

\begin{proposition}[La carga se conserva en cada etapa]\label{prop:g12:curly-arrows:charge}
La carga eléctrica total de las especies es la misma antes y después de
cada etapa de un mecanismo.
\end{proposition}

\begin{proof}
Una \omterm{def:g12:curly-arrows:curly-arrow}{flecha curva} solo desplaza \omterm{def:g9:inside-the-atom:nucleus}{electrones} que ya están presentes; no se
crea ni se destruye ningún \omterm{def:g9:inside-the-atom:nucleus}{electrón}, y los \omterm{def:g9:inside-the-atom:nucleus}{núcleos} no cambian.
\end{proof}

\begin{example}[Una primera flecha]\label{ex:g12:curly-arrows:first}
Un \omterm{def:g9:ions:ion}{ion} hidrógeno \ce{H+} se encuentra con un \omterm{def:g9:ions:ion}{ion} hidróxido \ce{OH-}. Un
\omterm{def:g10:lewis-and-shape:lone-pair}{par libre} del oxígeno forma un enlace con el \omterm{def:g9:ions:ion}{ion} hidrógeno, que no tiene
\omterm{def:g9:inside-the-atom:nucleus}{electrones}: \ce{H+ + OH- -> H2O}. Una sola flecha, desde el \omterm{def:g10:lewis-and-shape:lone-pair}{par libre} del
oxígeno hasta el \ce{H+}; las cargas $+1$ y $-1$ pasan a ser 0 en la
\omterm{def:g7:atoms-and-molecules:molecule}{molécula} de agua.
\end{example}

\begin{omfigure}
\setchemfig{atom sep=2.1em}
\chemfig{@{hp}H^{+}}\qquad\raisebox{0.2ex}{$+$}\qquad
\chemfig{H-@{ox}\charge{90=\:,0=\:,270=\:}{O}^{-}}
\qquad\raisebox{0.2ex}{$\longrightarrow$}\qquad
\chemfig{H-\charge{90=\:,270=\:}{O}-H}
\chemmove{
  \draw[->, thick, omProp, shorten <=4pt, shorten >=3pt]
    ($(ox)+(0,0.25)$) .. controls +(110:9mm) and +(60:9mm) .. (hp);}

{\small La etapa más sencilla: un \omterm{def:g10:lewis-and-shape:lone-pair}{par libre} del \omterm{def:g9:ions:ion}{ion} hidróxido se desplaza
para formar el nuevo enlace \ce{O-H} con el \omterm{def:g9:ions:ion}{ion} hidrógeno.}
\end{omfigure}

\section{Etapas elementales e intermedios}

\begin{definition}[Mecanismo, etapa elemental, intermedio]\label{def:g12:curly-arrows:mechanism}
El \emph{mecanismo de reacción}\index{mecanismo de reaccion@mecanismo de reacción} de una
reacción es la serie de etapas por las que los \omterm{def:g7:chemical-reactions:reactant}{reactivos} se convierten en
los productos. Cada \emph{etapa elemental}\index{etapa elemental} es un
único suceso en el que se desplazan a la vez unos pocos pares de
\omterm{def:g9:inside-the-atom:nucleus}{electrones}. Una especie formada en una etapa y consumida en otra
posterior es un \emph{intermedio de reacción}\index{intermedio de reaccion@intermedio de reacción}:
no aparece en la ecuación global.
\end{definition}


\begin{example}[Un carbocatión como intermedio]\label{ex:g12:curly-arrows:carbocation}
Cuando el bromuro de hidrógeno se adiciona al eteno, la reacción tiene
lugar en dos etapas. En la primera, el \omterm{def:g10:lewis-and-shape:multiple-bond}{doble enlace} del eteno cede un par
al hidrógeno del \ce{HBr}, y el par \ce{H-Br} va al bromo, que sale como
\ce{Br-}; el carbono que ha perdido su parte del \omterm{def:g10:lewis-and-shape:multiple-bond}{doble enlace} se queda
con solo seis \omterm{def:g9:inside-the-atom:nucleus}{electrones} y una carga positiva: un carbocatión,
\ce{CH3-CH2+}. En la segunda etapa, un \omterm{def:g10:lewis-and-shape:lone-pair}{par libre} del \ce{Br-} forma un
enlace con ese carbono. El carbocatión es un intermedio.
\end{example}

\section{Tres categorías de reacciones}

\begin{definition}[Sustitución, adición, eliminación]\label{def:g12:curly-arrows:categories}
En una \emph{sustitución}\index{sustitucion@sustitución}, un \omterm{def:g7:atoms-and-molecules:atom}{átomo} o un grupo de una
\omterm{def:g7:atoms-and-molecules:molecule}{molécula} se reemplaza por otro. En una \emph{adición}\index{adicion@adición}, dos
\omterm{def:g7:atoms-and-molecules:molecule}{moléculas} se unen en una sola, y un enlace múltiple pasa a ser simple. En
una \emph{eliminación}\index{eliminacion@eliminación}, se retira una \omterm{def:g7:atoms-and-molecules:molecule}{molécula} pequeña
(a menudo agua) de una \omterm{def:g7:atoms-and-molecules:molecule}{molécula}, y un enlace simple pasa a ser múltiple.
\end{definition}

\begin{omfigure}
\setchemfig{atom sep=2em}
\small
\textbf{\omterm{def:g12:curly-arrows:categories}{Sustitución}} (una etapa): \ce{CH3-CH2-Br + OH- -> CH3-CH2-OH + Br-}

\medskip
\chemfig{H-@{o1}\charge{90=\:,0=\:,270=\:}{O}^{-}}\qquad
\chemfig{H_3C-@{c1}C(-[2]H)(-[6]H)-[@{b1}]@{br1}\charge{90=\:,0=\:,270=\:}{Br}}
\qquad$\longrightarrow$\qquad
\chemfig{H_3C-C(-[2]H)(-[6]H)-O-H}\quad$+$\quad\ce{Br-}
\chemmove{
  \draw[->, thick, omProp, shorten <=4pt, shorten >=3pt]
    ($(o1)+(0.15,0.2)$) .. controls +(60:9mm) and +(120:9mm) .. (c1);
  \draw[->, thick, omDef, shorten <=2pt, shorten >=4pt]
    (b1) .. controls +(-90:7mm) and +(-90:7mm) .. (br1);}

\vspace{6mm}
\textbf{\omterm{def:g12:curly-arrows:categories}{Adición}} (dos etapas): \ce{CH2=CH2 + HBr -> CH3-CH2-Br}

\vspace{9mm}
\chemfig{H_2C=[@{pi}]CH_2}\quad$+$\quad
\chemfig{@{h2}H-[@{b2}]@{br2}Br}
\qquad$\longrightarrow$\qquad
\chemfig{H_3C-@{cp}\charge{90:2pt=$\scriptstyle +$}{C}H_2}\quad$+$\quad
\chemfig{@{brm}\charge{90=\:,180=\:,270=\:,0=\:}{Br}^{-}}
\qquad$\longrightarrow$\qquad \ce{CH3-CH2-Br}
\chemmove{
  \draw[->, thick, omProp, shorten <=2pt, shorten >=3pt]
    (pi) .. controls +(90:10mm) and +(110:9mm) .. (h2);
  \draw[->, thick, omDef, shorten <=2pt, shorten >=4pt]
    (b2) .. controls +(-90:7mm) and +(-90:7mm) .. (br2);
  \draw[->, thick, omProp, shorten <=4pt, shorten >=3pt]
    ($(brm)+(0,-0.25)$) .. controls +(-120:8mm) and +(-60:8mm) .. (cp);}

\vspace{6mm}
\textbf{\omterm{def:g12:curly-arrows:categories}{Eliminación}} (deshidratación ácida del etanol, tres etapas):
\ce{CH3-CH2-OH -> CH2=CH2 + H2O}

\medskip
\begin{tabular}{@{}l@{}}
1.\ el oxígeno capta un \omterm{def:g9:ions:ion}{ion} hidrógeno: \ce{CH3-CH2-OH + H+ -> CH3-CH2-OH2+};\\
2.\ el par \ce{C-O} sale con el agua: \ce{CH3-CH2-OH2+ -> CH3-CH2+ + H2O};\\
3.\ un par \ce{C-H} vecino pasa a ser el \omterm{def:g10:lewis-and-shape:multiple-bond}{doble enlace}, y el
    \ce{H+} sale:\\
\phantom{3.\ }\ce{CH3-CH2+ -> CH2=CH2 + H+}.
\end{tabular}

\medskip
{\small Tres mecanismos. Flechas naranjas: un par dador forma un nuevo
enlace. Flechas azules: un enlace se rompe y su par va al \omterm{def:g7:atoms-and-molecules:atom}{átomo} más
electronegativo. En la \omterm{def:g12:curly-arrows:categories}{eliminación}, el \omterm{def:g9:ions:ion}{ion} hidrógeno se recupera al
final.}
\end{omfigure}


\begin{remark}[Cambiar el esqueleto o el grupo]\label{rem:g12:curly-arrows:skeleton}
Algunas reacciones solo cambian el \omterm{def:g11:functional-groups:functional-group}{grupo funcional} y conservan el
esqueleto carbonado: la \omterm{def:g12:curly-arrows:categories}{sustitución} del bromoetano convierte un
\omterm{def:g11:functional-groups:halogenoalkane}{halogenoalcano} en un \omterm{def:g11:functional-groups:alcohol}{alcohol}, con dos carbonos antes y dos después. Otras
construyen o cortan el esqueleto: formar un nuevo enlace carbono--carbono
alarga una cadena, la etapa clave para fabricar una \omterm{def:g7:atoms-and-molecules:molecule}{molécula} grande a
partir de otras pequeñas. Planificar una \omterm{def:g10:synthesis-yield:synthesis}{síntesis} consiste en elegir
reacciones de ambos tipos, en el orden adecuado.
\end{remark}

\section{Ejercicios}

\begin{exercise}[$\star$]\label{exo:g12:curly-arrows:1}
En cada especie, nombra un sitio dador y, si lo hay, un sitio aceptor:
\ce{OH-}; \ce{H+}; \ce{CH2=CH2}; \ce{CH3-Cl}; \ce{H2O}.
\end{exercise}

\begin{exercise}[$\star$]\label{exo:g12:curly-arrows:2}
¿\omterm{def:g12:curly-arrows:categories}{Sustitución}, \omterm{def:g12:curly-arrows:categories}{adición} o \omterm{def:g12:curly-arrows:categories}{eliminación}?
(a) \ce{CH3-CH2-Cl + OH- -> CH3-CH2-OH + Cl-};
(b) \ce{CH2=CH2 + H2O -> CH3-CH2-OH};
(c) \ce{CH3-CH2-OH -> CH2=CH2 + H2O};
(d) \ce{CH2=CH2 + Br2 -> CH2Br-CH2Br}.
\end{exercise}

\begin{exercise}[$\star$]\label{exo:g12:curly-arrows:3}
¿Dónde empieza una \omterm{def:g12:curly-arrows:curly-arrow}{flecha curva} y dónde termina? ¿Qué significa una
\omterm{def:g12:curly-arrows:curly-arrow}{flecha curva} que va de un \omterm{def:g10:lewis-and-shape:lone-pair}{par libre} de un \omterm{def:g7:atoms-and-molecules:atom}{átomo} de oxígeno a un \omterm{def:g7:atoms-and-molecules:atom}{átomo} de
carbono?
\end{exercise}

\begin{exercise}[$\star$]\label{exo:g12:curly-arrows:4}
En la reacción \ce{H+ + NH3 -> NH4+}, ¿qué especie cede el par? Dibuja la
flecha. ¿Cuál es la carga del producto?
\end{exercise}

\begin{exercise}[$\star$]\label{exo:g12:curly-arrows:5}
¿Qué es un \omterm{def:g12:curly-arrows:mechanism}{intermedio de reacción}? Nombra el intermedio de la \omterm{def:g12:curly-arrows:categories}{adición} de
\ce{HBr} al eteno.
\end{exercise}

\begin{exercise}[$\star\star$]\label{exo:g12:curly-arrows:6}
Dibuja las \omterm{def:g12:curly-arrows:curly-arrow}{flechas curvas} de la \omterm{def:g12:curly-arrows:categories}{sustitución}
\ce{CH3-Br + OH- -> CH3-OH + Br-} (una etapa) y comprueba las cargas.
\end{exercise}

\begin{exercise}[$\star\star$]\label{exo:g12:curly-arrows:7}
En la primera etapa de la \omterm{def:g12:curly-arrows:categories}{adición} de \ce{HBr} al eteno, ¿por qué el par
\ce{H-Br} va al bromo y no al hidrógeno? ¿Qué carga lleva cada producto
de la etapa?
\end{exercise}

\begin{exercise}[$\star\star$]\label{exo:g12:curly-arrows:8}
En la deshidratación del etanol, escribe las tres etapas e indica, para
cada una, qué especie es el dador y cuál el aceptor. ¿Por qué no se
escribe el \omterm{def:g9:ions:ion}{ion} hidrógeno en la ecuación global?
\end{exercise}

\begin{exercise}[$\star\star$]\label{exo:g12:curly-arrows:9}
La etapa \ce{CH3-CH2+ + H2O -> CH3-CH2-OH2+} sigue a la \omterm{def:g12:curly-arrows:categories}{adición} de un
\omterm{def:g9:ions:ion}{ion} hidrógeno al eteno en agua. Dibuja la flecha. ¿Qué se forma tras una
nueva pérdida de \ce{H+}? Escribe la ecuación global e indica su
categoría.
\end{exercise}

\begin{exercise}[$\star\star$]\label{exo:g12:curly-arrows:10}
Lee el mecanismo de la \omterm{def:g12:curly-arrows:categories}{sustitución} del bromoetano: ¿cuántos pares se
desplazan? ¿Qué enlace se forma y cuál se rompe? ¿Hay un intermedio?
\end{exercise}

\begin{exercise}[$\star\star$]\label{exo:g12:curly-arrows:11}
% ledger: en:C, en:Cl, en:Br
El clorometano reacciona con el \omterm{def:g9:ions:ion}{ion} hidróxido más despacio que el
bromometano. Con las \omterm{def:g11:polarity-and-cohesion:electronegativity}{electronegatividades} (C 2.55, Cl 3.16, Br 2.96),
explica qué carbono tiene más carácter $\delta^+$. ¿Explica eso las
velocidades? (Piensa también en la facilidad con la que el \omterm{def:g9:ions:ion}{ion} saliente
se queda con el par.)
\end{exercise}

\begin{exercise}[$\star\star\star$]\label{exo:g12:curly-arrows:12}
El cloruro de hidrógeno se adiciona al propeno, \ce{CH2=CH-CH3}. En la
primera etapa, el \omterm{def:g9:ions:ion}{ion} hidrógeno puede unirse a cualquiera de los dos
carbonos del \omterm{def:g10:lewis-and-shape:multiple-bond}{doble enlace}, lo que da dos carbocationes posibles.
Dibújalos, junto con los dos productos posibles, y nómbralos. (Cuál se
forma mayoritariamente se estudia en el volumen del primer año.)
\end{exercise}

\begin{exercise}[$\star\star\star$]\label{exo:g12:curly-arrows:13}
Un alumno dibuja la primera etapa de la \omterm{def:g12:curly-arrows:categories}{sustitución} del bromoetano con
una flecha que va del \omterm{def:g7:atoms-and-molecules:atom}{átomo} de bromo al oxígeno del \ce{OH-}. Explica los
dos errores.
\end{exercise}

\begin{exercise}[$\star\star\star$]\label{exo:g12:curly-arrows:14}
En la formación de un \omterm{def:g11:functional-groups:carboxylic-acid}{éster} a partir de un \omterm{def:g11:functional-groups:carboxylic-acid}{ácido carboxílico} y un
\omterm{def:g11:functional-groups:alcohol}{alcohol}, el oxígeno del \omterm{def:g11:functional-groups:alcohol}{alcohol} ataca al carbono del grupo \ce{C=O}. Con
las \omterm{def:g11:polarity-and-cohesion:polar-bond}{cargas parciales}, explica por qué ese carbono es un sitio aceptor y
ese oxígeno un sitio dador. ¿Qué \omterm{def:g7:atoms-and-molecules:molecule}{molécula} pequeña se pierde en conjunto,
y a qué categoría pertenece la reacción completa?
\end{exercise}

\begin{exercise}[$\star\star\star$]\label{exo:g12:curly-arrows:15}
El eteno reacciona con el bromo \ce{Br2}, una \omterm{def:g11:polarity-and-cohesion:polar-molecule}{molécula apolar}. Cuando una
\omterm{def:g7:atoms-and-molecules:molecule}{molécula} de \ce{Br2} se acerca al \omterm{def:g10:lewis-and-shape:multiple-bond}{doble enlace}, los \omterm{def:g9:inside-the-atom:nucleus}{electrones} del
\ce{Br-Br} son empujados lejos del \omterm{def:g7:atoms-and-molecules:atom}{átomo} de bromo más cercano. Explica por
qué esto crea un sitio aceptor y propón la primera etapa del mecanismo.
\end{exercise}

\section{Problema: fabricar bromoetano}

\begin{problem}[{Problema de fin de semana --- a partir de diez gramos de
eteno: ¿cuánto bromoetano?}]
\label{pb:g12:curly-arrows:1}
% ledger: aw:C, aw:H, aw:Br, en:H, en:Br, en:C
El bromoetano, material de partida de muchas \omterm{def:g10:synthesis-yield:synthesis}{síntesis}, se fabrica
adicionando bromuro de hidrógeno al eteno: \ce{CH2=CH2 + HBr -> CH3-CH2-Br}.
Un químico usa \qty{10.0}{g} de eteno y un exceso de bromuro de
hidrógeno; el \omterm{def:g10:synthesis-yield:yield}{rendimiento} es del \qty{75}{\percent}.

\medskip
\noindent\textbf{Parte I --- Los sitios.}
\begin{enumerate}
  \item Dibuja las \omterm{def:g10:lewis-and-shape:lewis-structure}{estructuras de Lewis} del eteno y del bromuro de
        hidrógeno.
  \item ¿Dónde está el sitio dador del eteno? ¿Por qué?
  \item Calcula la diferencia de \omterm{def:g11:polarity-and-cohesion:electronegativity}{electronegatividad} del enlace \ce{H-Br}
        (H 2.20, Br 2.96). ¿Qué \omterm{def:g7:atoms-and-molecules:atom}{átomo} lleva $\delta^+$?
  \item ¿Qué \omterm{def:g7:atoms-and-molecules:atom}{átomo} del \ce{HBr} es el sitio aceptor?
\end{enumerate}

\medskip
\noindent\textbf{Parte II --- El mecanismo.}
\begin{enumerate}[resume]
  \item Dibuja la primera etapa con sus dos \omterm{def:g12:curly-arrows:curly-arrow}{flechas curvas}.
  \item Da las dos especies formadas y sus cargas.
  \item Comprueba la conservación de la carga en la primera etapa.
  \item Dibuja la segunda etapa con su \omterm{def:g12:curly-arrows:curly-arrow}{flecha curva}.
  \item ¿Cuántos \omterm{def:g9:inside-the-atom:nucleus}{electrones} rodean al carbono positivo del intermedio?
        ¿Por qué es tan reactivo?
\end{enumerate}

\medskip
\noindent\textbf{Parte III --- ¿Qué tipo de reacción?}
\begin{enumerate}[resume]
  \item ¿Cuál es el intermedio del mecanismo?
  \item ¿Por qué no aparece en la ecuación global?
  \item ¿A qué categoría pertenece la reacción?
  \item ¿Cambia la reacción el esqueleto carbonado, o solo el \omterm{def:g11:functional-groups:functional-group}{grupo
        funcional}?
\end{enumerate}

\medskip
\noindent\textbf{Parte IV --- ¿Cuánto producto?}
\begin{enumerate}[resume]
  \item Calcula las \omterm{def:g10:the-mole:molar-mass}{masas molares} del eteno y del bromoetano.
  \item Calcula la cantidad de eteno.
  \item Completa la \omterm{def:g10:reaction-progress-table:extent}{tabla de avance}. ¿Qué \omterm{def:g7:chemical-reactions:reactant}{reactivo} es limitante?
  \item Calcula la masa máxima de bromoetano.
  \item Calcula la masa obtenida con un \omterm{def:g10:synthesis-yield:yield}{rendimiento} del \qty{75}{\percent}.
  \item Da la respuesta final: ¿qué masa de bromoetano obtiene el químico?
\end{enumerate}
\end{problem}
